\heading{数据结构与算法A-17秋-期末}

oteinfo{本卷为 Python 语言版；题面与解答已逐题复核，图示均已改为随文 TikZ，不依赖外部图片文件。}

\begin{question}
一、选择填空题（7 小题，每空1 分，共10 分）\\
1.\\
下列说法中正确的是 \_\_\_\_\_\_\_\_\_\_\_\_。\\
A. 图的广度优先搜索中邻接点的寻找具有“先进后出”的特征。\\
B. 连通图的生成树是该图的一个极大连通子图。\\
C. 图的广度优先搜索是一个递归过程。\\
D. 在非连通图的遍历过程中，每调用一次深度优先搜索算法都得到该图的一\\
个连通分量。\\
2.\\
基于比较的内排序算法中，平均时间复杂度为O(nlogn)的算法 有\\
\_\_\_\_\_\_\_\_\_\_\_\_\_\_\_\_\_\_\_\_\_\_\_\_\_ ， 这其中稳定的排序算法 有\\
\_\_\_\_\_\_\_\_\_\_\_\_\_\_\_\_\_\_\_\_\_\_\_\_\_  。\\
3.\\
在包含n 个关键码的线性表中从后向前进行顺序检索（0 下标为监视哨）。若\\
第i 个关键码的检索概率为pi ，且满足如下分布 ：\\
p1 = 1/2, p2 = 1/4,  ..., pn = 1/2\\
n,\\
则成功检索的平均检索长度为  \_\_\_\_\_\_\_\_\_\_\_\_\_\_\_\_\_\_\_\_\_\_\_， 失败检索的平\\
均检索长度为 \_\_\_\_\_\_\_\_\_\_\_\_\_\_\_\_\_\_\_\_\_\_\_ 。\\
4.\\
给定一组输入关键码 \{15,  1,  3,  14,  11,  2,  12,  10,  4,  17\}，\\
采用大小为3 的最小堆进行置换选择排序的输出为\\
\_\_\_\_\_\_\_\_\_\_\_\_\_\_\_\_\_\_\_\_\_\_\_\_\_\_\_\_\_  。\\
5.\\
按照AVL 定义，下图所示的BST 中没有达到平衡状态的结点有\\
\_\_\_\_\_\_\_\_\_\_\_\_ 。\\
3 / 6\\
6.\\
假设磁盘块大小是1024 字节，每个索引项需要8 个字节。 一个包含20000\\
条记录的数据文件，若采用二级索引， 则一级索引文件和二级索引文件共需\\
占用 \_\_\_\_\_\_\_\_\_\_\_\_ 磁盘块。\\
7.\\
给定两个广义表 S = ((a，(b)，c)，((d)，e))， T = (f，(g，((h))，i，\\
j)); 利用广义表的Head 和Tail 运算， 则有\\
d = \_\_\_\_\_\_\_\_\_\_\_\_\_\_\_\_\_\_\_\_\_\_\_\_\_\_\_\_\_\_\_\_\_\_\_ ；\\
h = \_\_\_\_\_\_\_\_\_\_\_\_\_\_\_\_\_\_\_\_\_\_\_\_\_\_\_\_\_\_\_\_\_\_\_ 。\\
二、简答辨析题（5 小题，共15 分）\\
1. （2 分） 将关键码序列（7, 8, 30,11, 18, 9, 14）存储到一个散列表中。\\
散列表是一个下标从0 开始的一维数组， 若散列函数采用: H(key) = (key *\\
3) MOD 7， 处理冲突采用线性探测法，装填（负载）因子要求为0.7。\\
(1) . 请画出所构造的散列表;\\
(2) . 分别计算等概率情况下查找成功和查找不成功的平均查找长度。\\
2. （4 分）下图是一棵关键码为英文字符的m 阶B 树，依据字典序组织。 请回\\
答下述问题：\\
(1) . 这是一棵合法的B 树，则m 可以取哪些值，并说明原因;\\
(2) . 若m 取所有可能值中的最小值，则\\
(A) 给出查找字符V 的过程，并说明需要的读盘次数（假设根结点也需读\\
盘操作）；\\
(B) 给出插入字符I（H < I < J）的过程，并画出插入I 后的B 树；\\
(C) 在原始B 树（即进行第(B)问的操作之前）的基础上，给出删除H 的\\
过程，并画出删除H 后的B 树。\\
3. （4 分） 现有8 个顺串，每个顺串的第一个关键码为S1[1..8] = \{17, 10,\\
2, 9, 13, 15, 12, 4\}，第二个关键码是S2[1..8] =\{23, 34, 38, 25, 27,\\
40, 36, 39\}，进行从小到大的归并。\\
(1) . 请写出初始构建的赢者树内部结点的数组A[1..7] 和 败者树内部结点\\
4 / 6\\
的数组B[0..7]（数组A 和B 存储的是1 到8 的顺串的下标，B[0]表示最\\
终冠军）;\\
(2) . 在输出第一个全局赢者并重构赢者树或败者树后，写出此时赢者树内部\\
结点的数组A[1..7]和败者树内部结点的数组B[0..7]。\\
4. （3 分）请按照2, 1, 4, 5, 9 的插入顺序，建立一棵红黑树，并画出每插入\\
一个元素后的红黑树。（用B 和R 表示结点的颜色）\\
5. （2 分）给定下面的一棵splay 树， 按照结点序列2, 3, 4, 5, 6, 7 访问所\\
有结点之后，请画出访问后的splay 树形结构。\\
三、算法填空题（2 小题，每空1 分，共5 分）\\
1.\\
下面是考虑墓碑问题的散列表插入算法，算法中不允许有重复关键码插入。\\
template <class Key, class Elem, class KEComp, class EEComp> class\\
hashdict  \{\\
private:\\
Elem* HT;\\
// 散列表\\
int M;\\
// 散列表大小\\
int currcnt;\\
// 现有元素数目\\
Elem EMPTY;\\
// 空槽\\
int p(Key K, int i)\\
// 探查函数\\
int h(Key K) const ;\\
// 散列函数\\
Key getKey(Elem e);\\
// 获取元素e 的关键码值\\
public:\\
hashdict(int sz, Elem e)\{\\
// 构造函数，e 用来定义空槽\\
M = sz;\\
EMPTY = e;\\
currcnt = 0;\\
HT = Elem[sz]\\
for (int i=0; i<M; i++)\{\\
HT[i] = EMPTY;\\
\}\\
\}\\
\textasciitilde{}hashdict() \{ delete []HT; \}\\
5 / 6\\
hashSearch(const Key\&, Elem\&) const;\\
hashInsert(const Elem\&);\\
Elem hashDelete(const Key\& K);\\
int size() \{ return currcnt; \}\\
// 散列表中现有元素数\\
\};\\
// 墓碑用常量TOMB 表示，以下为插入函数\\
template<class Key,class Elem,class KEComp,class EEComp>\\
hashdict<Key,Elem,KEComp,EEComp>::hashInsert(const Elem\& e)\{\\
int insplace;\\
int i = 0;\\
int pos  = home = h(getkey(e));   // 找到新结点基地址\\
tomb\_pos = false;\\
while (\_\_\_\_\_ 填空1 \_\_\_\_\_\_\_\_\_)  \{\\
if (EEComp::eq(e,HT[pos]))\\
return False;\\
if (EEComp::eq(TOMB,HT[pos]) and !tomb\_pos) \{\\
\_\_\_\_\_\_\_\_\_ 填空2 \_\_\_\_\_\_\_\_\\
tomb\_pos = true;\\
\}\\
i++;\\
\_\_\_\_\_\_\_\_ 填空3 \_\_\_\_\_\_\_\_\_  // 探查\\
\}\\
if (!tomb\_pos)\\
insplace = pos;\\
// 若无墓碑，则insplace 位于空槽位置\\
HT[insplace] = e;\\
return True;\\
\}\\
2.\\
设有一个仅由红、白、蓝3 种颜色的条块组成的序列，各种色块的个数和分\\
布是随机的，但3 种颜色的总数为n。下面的代码实现了一种时间复杂度为\\
O(n)的排序，使得这些条块按照红、白、蓝的顺序排好。请将其补全。\\
Red, White, Blue = 1, 2, 3\\
def color\_sort(Array):\\
Insert\_Red, Insert\_White = 0, 0\\
for i in range(len(Array)):\\
if Array[i] == White:\\
Array[i] = Array[Insert\_White];\\
Array[Insert\_White] = White;\\
Insert\_White += 1\\
\}\\
if Array[i] == Red:\\
\_\_\_\_\_\_\_\_ 填空4 \_\_\_\_\_\_\_\_\_\\
\_\_\_\_\_\_\_\_ 填空5 \_\_\_\_\_\_\_\_\_\\
Array[Insert\_Red] = Red;\\
6 / 6\\
Insert\_Red += 1\\
Insert\_White += 1\\
\}\\
\}\\
\}\\
四、设计分析题（2 小题，共10 分）\\
1.\\
（6 分）请设计一个算法求有向无环图（DAG）中每对顶点间的“最长简单路\\
径”（所谓“最长简单路径”，是指该简单路径包含的边数最多）。以一个有向\\
无环图作为输入，对于每对顶点，如果它们之间存在简单路径，则输出其中路\\
径长度最长（边数最多）的简单路径，否则输出空。试分析你所设计算法的时\\
间复杂度。\\
2.\\
（4 分）利用证明基于比较的排序问题的复杂度下限的方法，试证明在已排序\\
序列上基于比较的检索问题的复杂度下限是Ω(logn)。\\
五、期中考试题（共30 分，成绩由助教登记）\\
六、上机考试题（共30 分，成绩由助教登记）\\
\\
\


\begin{solution}
一、选择与填空：
\begin{enumerate}[leftmargin=2em]
  \item 选 D。BFS 使用队列而非“先进后出”；生成树是极小连通子图而非极大连通子图；BFS 通常不是递归过程。非连通无向图中每次从未访问顶点调用 DFS 得到一个连通分量。
  \item 平均时间复杂度为 $O(n\log n)$ 的典型比较排序有快速排序、归并排序、堆排序；其中稳定的是归并排序。
  \item 若题意为 $p_i=2^{-i}$ 并按其归一化处理，则
  \[
  ASL_{succ}=\frac{\sum_{i=1}^n (n-i+1)2^{-i}}{1-2^{-n}}.
  \]
  从后向前顺序检索失败时需检查全部 $n$ 个关键码；若把 0 下标监视哨也计为一次比较，则 $ASL_{fail}=n+1$，否则为 $n$。
  \item 置换选择排序输出的初始归并段为：第一段 $1,3,11,14,15$；第二段 $2,4,10,12,17$。
  \item 该题原扫描中 BST 图像缺失；若按已补图版本作答，应逐个结点计算左右子树高度差，所有 $|BF|>1$ 的结点即为未平衡结点。
  \item 每个磁盘块可放 $1024/8=128$ 个索引项。一级索引需 $\lceil20000/128\rceil=157$ 块，二级索引需 $\lceil157/128\rceil=2$ 块，共 $159$ 块。
  \item $d=Head(Head(Head(Tail(S))))$；$h=Head(Head(Head(Tail(Head(Tail(T))))))$。
\end{enumerate}

二、简答辨析题：
\begin{enumerate}[leftmargin=2em]
  \item 散列表长度满足 $7/L\le0.7$，故取 $L=10$。依次插入 $7,8,30,11,18,9,14$，$H(key)=3key\bmod7$，线性探测得到：
  \[
  \begin{array}{c|cccccccccc}
  地址&0&1&2&3&4&5&6&7&8&9\\\hline
  关键码&7&14&--&8&--&11&30&18&9&--
  \end{array}
  \]
  成功比较次数为 $1,1,1,1,3,3,2$，$ASL_{succ}=12/7$。不成功查找从基地址 $0\sim6$ 出发的比较次数为 $3,2,1,2,1,5,4$，$ASL_{fail}=18/7$。
  \item B 树题依赖原图。作答原则：由每个内部结点关键字数确定 $m$ 的上下界；查找 $V$ 时自根向叶按关键字区间下降，访问一个结点算一次读盘；插入 $I$ 先定位叶结点，若溢出则分裂并向父结点上推中间关键字；删除 $H$ 若叶结点下溢，则优先向兄弟借关键字，否则合并并向父结点递归调整。
  \item 顺串归并题：第一轮全局最小来自第 3 个顺串，输出 $2$；随后第 8 个顺串首关键码 $4$ 成为下一全局最小。赢者树/败者树数组应按“存顺串下标而非关键码”的约定填写，重构时只更新输出顺串所在路径。
  \item 按 $2,1,4,5,9$ 插入红黑树：插入 $2$ 后根为黑；插入 $1$、$4$ 为红，树仍平衡；插入 $5$ 时父 $4$ 与叔 $1$ 同为红，变色为 $1B,4B,2B,5R$；插入 $9$ 触发以 $4$ 为祖父的 RR 调整，得到根 $2B$，左子 $1B$，右子 $5B$，$5$ 的左子 $4R$、右子 $9R$。
  \item Splay 题依赖原图。作答时应对访问序列 $2,3,4,5,6,7$ 逐次执行伸展：若访问结点父为根则单旋；若为 zig-zig 则先旋祖父再旋父；若为 zig-zag 则先旋父再旋祖父。最后一次访问的 $7$ 应位于根。
\end{enumerate}

三、算法填空题：
\begin{enumerate}[leftmargin=2em]
  \item 墓碑散列表插入：
  \begin{enumerate}[leftmargin=2em]
    \item \texttt{\detokenize{not EEComp.eq(EMPTY, HT[pos]) and i < M}}；
    \item \texttt{\detokenize{insplace = pos}}；
    \item \texttt{pos = (home + p(getKey(e), i)) \% M}。
  \end{enumerate}
  循环结束后，若此前没有遇到墓碑，则插入到第一个空槽；若遇到墓碑，则复用第一个墓碑位置。
  \item 三色排序填空：
  \begin{enumerate}[leftmargin=2em]
    \item \texttt{\detokenize{Array[i] = Array[Insert_White]}}；
    \item \texttt{\detokenize{Array[Insert_White] = Array[Insert_Red]}}。
  \end{enumerate}
  然后题中已有语句 \texttt{\detokenize{Array[Insert_Red] = Red}}，并同时推进两个插入边界。
\end{enumerate}

四、设计分析题：
\begin{enumerate}[leftmargin=2em]
  \item 对 DAG 先做拓扑排序。对每个源点 $s$，按拓扑序进行动态规划：\texttt{\detokenize{dist[s]=0}}，其余为 $-\infty$；扫描每条边 $(u,v)$，若 \texttt{\detokenize{dist[u]+1>dist[v]}} 则更新 \texttt{\detokenize{dist[v]}} 与前驱。对每个 $s$ 做一次即可得到从 $s$ 到所有点的最长简单路径。时间复杂度为 $O(|V|(|V|+|E|))$，若只求长度可直接保存距离，若要输出路径则同时保存前驱并回溯。
  \item 基于比较的有序表检索可看作决策树。若要判断目标属于哪个位置或哪个间隙，至少有 $n+1$ 个可能结果；一棵二叉比较决策树高为 $h$ 时最多区分 $2^h$ 个结果，故 $2^h\ge n+1$，即 $h\ge\lceil\log_2(n+1)\rceil=\Omega(\log n)$。
\end{enumerate}
\end{solution}

\end{question}
